% =====================================================================
%  OpenFlexure USB — Guide accessible à tous
%  Compilation : pdfLaTeX (Overleaf par défaut) ; XeLaTeX et LuaLaTeX acceptés.
% =====================================================================
\documentclass[11pt,a4paper,openany]{report}
% =====================================================================
%  Préambule — compatible pdfLaTeX (Overleaf par défaut), XeLaTeX et LuaLaTeX
% =====================================================================
\usepackage{iftex}
\ifPDFTeX
  \usepackage[T1]{fontenc}
  \usepackage[utf8]{inputenc}
\fi
\usepackage{amsmath}
\usepackage[french]{babel}
\usepackage{libertinus}
\usepackage{microtype}

\usepackage[a4paper,top=2.5cm,bottom=2.5cm,left=2.4cm,right=2.4cm,headheight=14pt]{geometry}
\usepackage[table]{xcolor}
\usepackage{graphicx}

\usepackage[locale=FR,per-mode=symbol,group-minimum-digits=4]{siunitx}
\sisetup{range-phrase={\,\text{à}\,},list-final-separator={ et },inter-unit-product=\ensuremath{\cdot}}

\usepackage{booktabs}
\usepackage{tabularx}
\usepackage{longtable}
\usepackage{multirow}
\usepackage{array}
\usepackage{makecell}
\usepackage{caption}
\usepackage{subcaption}
\usepackage{float}
\usepackage{enumitem}
\usepackage{pdflscape}
\usepackage{eso-pic}

\usepackage{pifont}
\usepackage{tikz}
\usetikzlibrary{babel}
\usetikzlibrary{arrows.meta,positioning,fit,calc}

\usepackage[most]{tcolorbox}
\usepackage{listings}
\usepackage{fancyhdr}
\usepackage{titlesec}
\usepackage[hidelinks]{hyperref}
\usepackage[french,nameinlink]{cleveref}

% ---------------------------------------------------------------------
%  Couleurs
% ---------------------------------------------------------------------
\definecolor{pslbleu}{HTML}{1F4E8C}
\definecolor{pslbleuclair}{HTML}{E8EEF7}
\definecolor{pslorange}{HTML}{E8781A}
\definecolor{pslorangeclair}{HTML}{FDF0E4}
\definecolor{pslvert}{HTML}{2E7D4F}
\definecolor{pslvertclair}{HTML}{E7F3EC}
\definecolor{pslrouge}{HTML}{B3261E}
\definecolor{pslrougeclair}{HTML}{FBEAE9}
\definecolor{pslgris}{HTML}{5F6368}

\hypersetup{colorlinks=true,linkcolor=pslbleu,urlcolor=pslbleu,citecolor=pslbleu,
  pdftitle={OpenFlexure USB — Microscope en périphérique USB},
  pdfauthor={Tom Penfornis},pdfsubject={Liaison USB gadget et mise à l'heure}}

% ---------------------------------------------------------------------
%  Titres et en-têtes
% ---------------------------------------------------------------------
\titleformat{\chapter}[display]
  {\normalfont\sffamily\color{pslbleu}}
  {\large\MakeUppercase{\chaptertitlename}~\thechapter}
  {4pt}
  {\Huge\bfseries}
  [\vspace{4pt}{\color{pslorange}\titlerule[1.5pt]}]
\titlespacing*{\chapter}{0pt}{-10pt}{24pt}
\titleformat{\section}{\normalfont\sffamily\Large\bfseries\color{pslbleu}}{\thesection}{0.8em}{}
\titleformat{\subsection}{\normalfont\sffamily\large\bfseries\color{pslbleu!85!black}}{\thesubsection}{0.8em}{}
\titleformat{\subsubsection}{\normalfont\sffamily\normalsize\bfseries}{\thesubsubsection}{0.8em}{}

\pagestyle{fancy}
\fancyhf{}
\fancyhead[L]{\small\sffamily\color{pslgris}OpenFlexure USB — Rapport technique}
\fancyhead[R]{\small\sffamily\color{pslgris}\nouppercase{\leftmark}}
\fancyfoot[C]{\small\sffamily\thepage}
\renewcommand{\headrulewidth}{0.4pt}
\renewcommand{\headrule}{{\color{pslbleu!40}\hrule height\headrulewidth width\headwidth\vskip-\headrulewidth}}
\fancypagestyle{plain}{\fancyhf{}\fancyfoot[C]{\small\sffamily\thepage}\renewcommand{\headrulewidth}{0pt}}

\setcounter{tocdepth}{2}
\setcounter{secnumdepth}{3}
\setlength{\parskip}{0.35em}
\setlist{itemsep=2pt,topsep=4pt}

\captionsetup{font=small,labelfont={bf,sf,color=pslbleu},format=hang,margin=10pt}
\captionsetup[table]{position=top}
\renewcommand{\arraystretch}{1.18}
\newcolumntype{L}{>{\raggedright\arraybackslash}X}
\newcolumntype{C}{>{\centering\arraybackslash}X}

% ---------------------------------------------------------------------
%  Encadrés
% ---------------------------------------------------------------------
\newtcolorbox{synthese}[1][Synthèse]{enhanced,breakable,colback=pslbleuclair,colframe=pslbleu,
  fonttitle=\sffamily\bfseries,title={#1},arc=2pt,boxrule=0.6pt,left=6pt,right=6pt}
\newtcolorbox{choix}[1][Choix retenu]{enhanced,breakable,colback=pslvertclair,colframe=pslvert,
  fonttitle=\sffamily\bfseries,title={#1},arc=2pt,boxrule=0.6pt,left=6pt,right=6pt}
\newtcolorbox{attention}[1][Point de vigilance]{enhanced,breakable,colback=pslrougeclair,colframe=pslrouge,
  fonttitle=\sffamily\bfseries,title={#1},arc=2pt,boxrule=0.6pt,left=6pt,right=6pt}
\newtcolorbox{calcul}[1][Calcul]{enhanced,breakable,colback=white,colframe=pslorange,
  fonttitle=\sffamily\bfseries,coltitle=black,colbacktitle=pslorangeclair,title={#1},arc=2pt,boxrule=0.6pt,left=6pt,right=6pt}

% ---------------------------------------------------------------------
%  Code
% ---------------------------------------------------------------------
\lstdefinestyle{psl}{basicstyle=\ttfamily\footnotesize,keywordstyle=\color{pslbleu},
  commentstyle=\color{pslgris}\itshape,stringstyle=\color{pslorange!80!black},numbers=left,
  numberstyle=\tiny\color{pslgris},frame=single,rulecolor=\color{pslbleu!30},backgroundcolor=\color{pslbleuclair!40},
  breaklines=true,columns=fullflexible,keepspaces=true,showstringspaces=false,
  literate={é}{{\'e}}1 {è}{{\`e}}1 {à}{{\`a}}1 {ç}{{\c{c}}}1 {ê}{{\^e}}1 {É}{{\'E}}1 {ô}{{\^o}}1 {î}{{\^i}}1 {û}{{\^u}}1 {ù}{{\`u}}1 {Ø}{{\O}}1 {×}{{$\times$}}1 {â}{{\^a}}1 {⇄}{{$\leftrightarrows$}}1 {→}{{$\rightarrow$}}1 {«}{{\guillemotleft}}1 {»}{{\guillemotright}}1 {À}{{\`A}}1}
\lstset{style=psl}

% ---------------------------------------------------------------------
%  Raccourcis
% ---------------------------------------------------------------------
\newcommand{\comp}[1]{\textsf{#1}}
\newcommand{\fichier}[1]{{\small\nolinkurl{#1}}}

% Pas de citation (lien dans un lien) dans les listes des figures et des tableaux
\let\listoffiguresorig\listoffigures
\let\listoftablesorig\listoftables
\renewcommand{\listoffigures}{\begingroup\renewcommand{\cite}[2][]{}\listoffiguresorig\endgroup}
\renewcommand{\listoftables}{\begingroup\renewcommand{\cite}[2][]{}\listoftablesorig\endgroup}
\crefname{lstlisting}{listing}{listings}
\Crefname{lstlisting}{Listing}{Listings}


\hypersetup{pdftitle={Un microscope OpenFlexure utilisable partout, sans réseau},
  pdfsubject={Guide d'utilisation de la liaison USB}}
\fancyhead[L]{\small\sffamily\color{pslgris}OpenFlexure USB — Guide}
\newcommand{\ip}[1]{\texttt{#1}}
\newcommand{\oui}{\textcolor{pslvert}{\ding{51}}}
\newcommand{\nonteste}{\textcolor{pslorange}{\ding{108}}}
\newcommand{\marque}[1]{\tikz[remember picture,overlay]\coordinate (#1);}

\begin{document}

\pagenumbering{Alph}
\begin{titlepage}
\begin{tikzpicture}[remember picture,overlay]
  \fill[pslbleu] (current page.north west) rectangle ([yshift=-8.2cm]current page.north east);
  \fill[pslorange] ([yshift=-8.2cm]current page.north west) rectangle ([yshift=-8.5cm]current page.north east);
  \node[anchor=north west,text=white!85!pslbleu,font=\sffamily\small,inner sep=0pt] at ([xshift=2.4cm,yshift=-1.7cm]current page.north west) {Projet OpenFlexure USB};
  \node[anchor=north west,text=white,font=\sffamily\bfseries\fontsize{30}{36}\selectfont,inner sep=0pt,text width=16.5cm] at ([xshift=2.4cm,yshift=-2.6cm]current page.north west) {Un microscope utilisable partout, sans réseau};
  \node[anchor=north west,text=white!90!pslbleu,font=\sffamily\Large,inner sep=0pt,text width=16cm] at ([xshift=2.4cm,yshift=-6.25cm]current page.north west) {Un câble USB, un ordinateur, un navigateur};
\end{tikzpicture}

\vspace*{8.2cm}
\begin{center}
\begin{tikzpicture}[font=\sffamily\small,
  boite/.style={draw=pslbleu,fill=pslbleuclair,rounded corners=3pt,minimum height=1.5cm,align=center,line width=0.8pt},
  lien/.style={line width=1.4pt,pslorange}]
  \node[boite,minimum width=3.8cm] (pc) {\textbf{Votre ordinateur}\\navigateur web};
  \node[boite,minimum width=3.8cm,right=4.2cm of pc] (pi) {\textbf{Microscope}\\OpenFlexure};
  \draw[lien,{Latex}-{Latex}] (pc) -- node[above,text=black]{câble USB-C} node[below,text=black,font=\sffamily\footnotesize]{ni Wi-Fi, ni Internet} (pi);
\end{tikzpicture}
\end{center}
\vfill
{\sffamily
\begin{tabularx}{\textwidth}{@{}lL@{}}
\toprule
\textbf{Pour qui} & Chercheurs sur le terrain, enseignants, makers et toute personne utilisant un microscope OpenFlexure sans réseau disponible \\
\textbf{Prérequis} & Aucune connaissance en informatique pour l'utilisation ; quelques commandes pour l'installation \\
\textbf{Auteur} & Tom Penfornis \\
\textbf{Version} & 1.0 --- 30 octobre 2026 \\
\bottomrule
\end{tabularx}}
\end{titlepage}

\pagenumbering{roman}
\chapter*{L'essentiel en une page}
\addcontentsline{toc}{chapter}{L'essentiel en une page}

\begin{synthese}[En bref]
Le microscope OpenFlexure est un microscope open source, fabriqué en grande partie par impression 3D. On l'utilise depuis un navigateur web, mais il faut pour cela le relier à un réseau.

Ce projet supprime cette contrainte : \textbf{on branche le microscope à un ordinateur avec un câble USB-C, on ouvre une page web, et c'est tout.}
\end{synthese}

\begin{itemize}
  \item \textbf{Sans connexion :} aucun Wi-Fi, Internet ou box n'est nécessaire.
  \item \textbf{Sans logiciel à installer} sur l'ordinateur : un navigateur suffit.
  \item \textbf{Images bien datées :} le microscope récupère l'heure de l'ordinateur, même hors ligne.
  \item \textbf{Un seul câble possible :} il peut transporter à la fois l'alimentation et les commandes.
  \item \textbf{Réversible :} l'installation peut être désinstallée.
\end{itemize}

\begin{attention}[Où en est le projet ?]
La liaison fonctionne sous Linux, Windows et macOS. En revanche, \textbf{le montage n'a pas encore été utilisé sur un vrai terrain}. Le \cref{chap-etat} détaille ce qui est prouvé et ce qui ne l'est pas.
\end{attention}

\tableofcontents
\cleardoublepage
\pagenumbering{arabic}

% =====================================================================
\chapter{Pourquoi ce projet ?}
% =====================================================================

\section{Le problème}
Le microscope OpenFlexure contient un petit ordinateur, un Raspberry Pi. Pour le commander depuis un autre ordinateur, on passe habituellement par un réseau. Or, ce réseau n'existe pas toujours :
\begin{itemize}
  \item \textbf{sur le terrain} (mission, station isolée, zone sans couverture), il n'y a parfois ni connexion ni box Wi-Fi ;
  \item \textbf{dans une école ou un petit laboratoire}, l'accès au réseau peut être indisponible ou compliqué à obtenir ;
  \item \textbf{dans certains établissements}, le réseau est réservé aux appareils enregistrés.
\end{itemize}

\section{La solution}
Le microscope se présente à l'ordinateur comme une \textbf{petite carte réseau USB}. Il suffit d'un câble entre les deux. C'est comme un mini-réseau privé de deux appareils : l'ordinateur et le microscope, rien d'autre.

\begin{figure}[H]
\centering
\begin{tikzpicture}[font=\sffamily\small,
  boite/.style={draw=pslbleu,fill=pslbleuclair,rounded corners=3pt,minimum height=1.2cm,minimum width=3.2cm,align=center},
  barre/.style={draw=pslgris,fill=white,rounded corners=3pt,minimum height=1.2cm,minimum width=3.2cm,align=center,text=pslgris},
  fl/.style={{Latex}-{Latex},line width=1.2pt}]
  \node[font=\sffamily\bfseries,text=pslgris] at (3.5,1.2) {Avant};
  \node[boite] (a1) at (0,0) {Ordinateur};
  \node[barre] (a2) at (3.5,0) {Réseau\\obligatoire};
  \node[boite] (a3) at (7,0) {Microscope};
  \draw[fl,pslgris] (a1) -- (a2);
  \draw[fl,pslgris] (a2) -- (a3);
  \node[font=\sffamily\bfseries,text=pslorange!80!black] at (3.5,-1.3) {Avec ce projet};
  \node[boite] (b1) at (0,-2.5) {Ordinateur};
  \node[boite] (b3) at (7,-2.5) {Microscope};
  \draw[fl,pslorange] (b1) -- node[above,text=black]{câble USB-C} (b3);
\end{tikzpicture}
\caption{Le réseau intermédiaire est remplacé par un simple câble}
\end{figure}

\section{Et l'heure des images ?}
Une image scientifique doit être correctement datée. Or, le Raspberry Pi n'a pas d'horloge à pile : sans Internet, il repart de la dernière heure connue, qui peut dater de plusieurs jours.

La page d'accueil règle ce problème : quand vous l'ouvrez, votre ordinateur transmet son heure au microscope. Il n'y a rien à faire, mais il faut bien passer par cette page. La date des images dépend donc de la bonne heure de votre ordinateur.

% =====================================================================
\chapter{Utiliser le microscope}
% =====================================================================

\section{Ce qu'il vous faut}
\begin{itemize}
  \item un microscope OpenFlexure équipé d'un Raspberry Pi~4 (4\,Go ou plus) et d'une Sangaboard~v5, avec la liaison USB installée (\cref{chap-installation}) ;
  \item un ordinateur avec un navigateur web, de préférence avec un port USB-C ;
  \item un câble USB-C capable de transporter des \textbf{données} : certains câbles ne servent qu'à la recharge.
\end{itemize}

\section{Le branchement : à lire en premier}
\begin{attention}[Une seule source de courant à la fois]
\textbf{Ne branchez jamais l'ordinateur au Raspberry Pi pendant que la Sangaboard reçoit du courant par son propre port USB-C}, sauf avec un bloqueur de 5\,V. Les deux sources seraient reliées directement, ce qui peut endommager le matériel.
\end{attention}

\begin{table}[H]
\centering
\caption{Les deux montages possibles}
\small
\begin{tabularx}{\textwidth}{@{}>{\bfseries}p{3cm}LL@{}}
\toprule
& \marque{cadrehaut}Un seul câble & \marque{cadredroite}Alimentation séparée \\
\midrule
Branchement & Ordinateur vers USB-C du Pi ; rien sur l'USB-C de la Sangaboard & Alimentation vers USB-C de la Sangaboard ; ordinateur vers USB-C du Pi à travers un bloqueur de 5\,V \\
Avantage & Très simple : idéal lorsqu'on n'a pas de prise électrique & Ne sollicite pas la batterie de l'ordinateur \\
Limite & Dépend de la puissance du port ; vide plus vite la batterie du portable & Demande une prise (ou une batterie externe) et un accessoire \\
État & Testé sur un seul ordinateur, moteurs compris\marque{cadrebas} & Pas encore testé \\
\bottomrule
\end{tabularx}
\begin{tikzpicture}[remember picture,overlay]
  \draw[pslvert,line width=0.5pt,rounded corners=2pt] ($(cadrehaut)+(-5pt,12pt)$) rectangle ($(cadredroite |- cadrebas)+(-7pt,-5pt)$);
\end{tikzpicture}
\end{table}

\section{Pas à pas}
\begin{enumerate}
  \item Branchez le microscope selon l'un des deux montages.
  \item Attendez environ une minute, le temps que le microscope démarre.
  \item Sur l'ordinateur, ouvrez votre navigateur et saisissez \textbf{\ip{http://10.55.0.1}}. Enregistrez cette adresse en favori.
  \item Un message confirme la mise à l'heure, puis l'interface du microscope s'affiche.
  \item En fin de séance, cliquez sur le bouton « Éteindre » de l'interface, attendez l'arrêt, puis débranchez.
\end{enumerate}

\begin{choix}[Bon réflexe avant de partir sur le terrain]
Faites un essai complet à la maison ou au laboratoire, \textbf{Wi-Fi coupé} et sur batterie, avec l'ordinateur que vous emporterez. Vérifiez que les moteurs fonctionnent, que l'heure est juste et notez l'autonomie obtenue. Le bouton d'extinction évite de corrompre la carte SD.
\end{choix}

\section{En cas de problème}
\begin{table}[H]
\centering
\caption{Problèmes courants}
\small
\begin{tabularx}{\textwidth}{@{}p{5.2cm}L@{}}
\toprule
Ce que vous constatez & Ce qu'il faut essayer \\
\midrule
La page ne s'ouvre pas & Attendez une minute après la mise sous tension, puis changez de câble : il ne transporte peut-être pas les données \\
Le microscope s'éteint ou redémarre & Le port USB ne fournit sans doute pas assez de courant : essayez un autre port USB-C ou une alimentation séparée \\
L'heure des images est fausse & Passez bien par \ip{http://10.55.0.1}, et non directement par l'adresse se terminant par \ip{:5000} ; vérifiez l'heure de l'ordinateur \\
Ça ne marche pas sur un ordinateur professionnel & Un proxy, un pare-feu ou une règle USB peut bloquer la liaison : demandez de l'aide au service informatique \\
\bottomrule
\end{tabularx}
\end{table}

% =====================================================================
\chapter{Où en est le projet ?}
\label{chap-etat}
% =====================================================================

Le tableau distingue ce qui a réellement été vérifié de ce qui reste à prouver.

\begin{table}[H]
\centering
\caption{État des tests (\oui{} vérifié, \nonteste{} pas encore vérifié)}
\small
\begin{tabularx}{\textwidth}{@{}LLc@{}}
\toprule
Point vérifié & Résultat & \\
\midrule
Connexion USB automatique & La liaison revient seule après un redémarrage & \oui \\
Ordinateur sous Linux & Fonctionne (Debian~12) & \oui \\
Ordinateur sous Windows & Fonctionne & \oui \\
Ordinateur sous macOS & Fonctionne & \oui \\
Mise à l'heure sans Internet & Microscope déréglé remis à l'heure, Wi-Fi coupé & \oui \\
Image copiée sur une carte SD neuve & Démarrage et fonctionnement validés & \oui \\
Moteurs alimentés par un seul câble & Fonctionne, sur un seul ordinateur & \oui \\
Autres ordinateurs portables & Le courant disponible varie selon le modèle & \nonteste \\
Autonomie sur batterie & Pas encore mesurée & \nonteste \\
Bloqueur de 5\,V & Pas encore essayé & \nonteste \\
Mission réelle (chaleur, poussière, transport) & Pas encore réalisée & \nonteste \\
\bottomrule
\end{tabularx}
\end{table}

Le matériel utilisé pour ces essais est un Raspberry Pi~4 avec une Sangaboard~v5 et la version bêta 3.0.0 d'OpenFlexure. Cette version est encore en développement : une mise à jour peut obliger à relancer l'installation.

Une batterie intégrée au boîtier est à l'étude. Elle n'est pas encore assez avancée pour être partagée.

% =====================================================================
\chapter{Installer la liaison USB}
\label{chap-installation}
% =====================================================================

Cette étape n'est à faire qu'une fois. Elle demande quelques commandes ; une personne habituée à Linux pourra s'en charger en quelques minutes.

\begin{enumerate}
  \item Installez OpenFlexure en version 3 bêta en suivant les instructions officielles : \url{https://openflexure.org/projects/microscope/install}.
  \item Téléchargez ce projet et copiez le script \fichier{scripts/ofm-usb-install.sh} sur le Raspberry Pi.
  \item Sur le Raspberry Pi, lancez les commandes suivantes :
\end{enumerate}
\begin{lstlisting}[language=bash,numbers=none]
sudo ./ofm-usb-install.sh / /boot/firmware
sudo reboot
\end{lstlisting}

Pour tout retirer, utilisez la même commande en ajoutant \ip{-{}-uninstall} juste après le nom du script. Les détails techniques et la liste des éléments modifiés se trouvent dans le fichier \fichier{README.md} du projet.

\begin{attention}[Vous partagez une carte SD avec d'autres personnes ?]
Créez une image avec \fichier{scripts/ofm-image.sh} : elle retire les mots de passe, le Wi-Fi, les images enregistrées et l'étalonnage, puis liste les fichiers restants. Relisez cette liste avant de partager. Le fichier \fichier{README.md} explique la marche à suivre.
\end{attention}

% =====================================================================
\chapter{Questions fréquentes}
% =====================================================================

\paragraph{Ai-je besoin d'Internet, même une seule fois ?}
Pour utiliser le microscope, non. Internet ne sert qu'à l'installation initiale, pour télécharger les fichiers.

\paragraph{Mon ordinateur perd-il sa connexion quand le microscope est branché ?}
Non. Le câble ne sert qu'à parler au microscope ; l'ordinateur conserve sa connexion habituelle s'il en a une.

\paragraph{Le microscope est-il accessible à d'autres personnes ?}
Avec cette liaison, seul l'ordinateur relié par le câble peut le joindre. Si le Wi-Fi du microscope reste activé, il reste aussi joignable par ce réseau.

\paragraph{Peut-on l'utiliser dans un hôpital ou une entreprise ?}
Le microscope n'est pas raccordé au réseau de l'établissement et aucune protection n'est contournée. Il faut cependant respecter ses règles sur les périphériques USB : demandez l'avis du service informatique. Ce montage n'est pas un dispositif médical.

\paragraph{Puis-je utiliser une tablette ou un téléphone ?}
Cela n'a pas été testé.

\paragraph{Est-ce gratuit ?}
Oui. Les scripts sont libres, sous licence GPL version~3 ou ultérieure. OpenFlexure est un projet indépendant : son logiciel serveur est sous GPL~3 et son matériel sous licence CERN-OHL-S~2.0.

% =====================================================================
\chapter{Petit lexique}
% =====================================================================

\begin{table}[H]
\centering
\small
\begin{tabularx}{\textwidth}{@{}>{\bfseries}p{3.8cm}L@{}}
\toprule
Bloqueur de 5\,V & Petit adaptateur USB qui laisse passer les données mais pas le courant. Attention, il ne s'agit pas d'un « data blocker », qui fait l'inverse \\
Carte SD & Petite carte mémoire contenant le système du microscope \\
Navigateur & Logiciel pour afficher des pages web (Firefox, Chrome, Edge\dots) \\
Raspberry Pi & Petit ordinateur intégré au microscope \\
Sangaboard & Carte électronique qui commande les moteurs \\
USB-C & Prise USB moderne, ovale et réversible \\
\ip{10.55.0.1} & Adresse du microscope sur le mini-réseau créé par le câble \\
\bottomrule
\end{tabularx}
\end{table}

\end{document}
